\section{Conclusion}
\label{sec:concl}

This paper presented RPL-ClusterIDS, a clustering-based, energy-aware, and context-adaptive hierarchical distributed intrusion detection system for resource-constrained RPL networks.
The key idea is to treat clustering not as a passive grouping structure but as the mechanism that allocates detection depth, communication budget, and cluster-head responsibility according to residual energy, topological stability, and traffic priority.
Members perform ultra-light screening through four rules and a compact local anomaly score.
Cluster heads aggregate evidence, apply advanced rules, and run a two-stage threshold verification for confirmation.
A three-mode policy makes the security--lifetime trade-off explicit.

The Contiki-NG implementation and reproducibility package provide a foundation for Computer Networks submission. The pilot campaign (21 seeds for CLUSTERIDS, 3 seeds for the ablation study) indicates RPL-ClusterIDS achieves 80.6\% campaign-level detection rate with 0.50\% false-positive rate and 5--7\% CPU overhead, while the B1 baseline yields 0\% detection rate in all configurations tested. Full statistical validation across additional seeds is left for a future campaign.

All materials are openly available at \url{https://github.com/madani-belacel/IDS-IOT}. Future work will target hardware validation, robust cluster-head election under adversarial pressure, online model adaptation, sensitivity analysis completion, and integration with standardized RPL security mechanisms.

\section*{Acknowledgment}
The author thanks colleagues at Abdelhamid Ibn Badis University of Mostaganem for feedback on earlier drafts of this manuscript.
